\documentclass[12pt]{article}
\usepackage[T1]{fontenc}
\usepackage[utf8]{inputenc}
\usepackage{lmodern}
\usepackage{textcomp}
\usepackage{enumitem}
\usepackage{geometry}
\geometry{margin=1in}
\begin{document}
\begin{center}
Answer Key - Sociology - Undergraduate Level Assessment 16 \
\small (Answers are ordered exactly as in the question paper.)
\end{center}

\section*{Multiple Choice}
\begin{enumerate}[label=\arabic*.]
\item \textbf{Answer:} B
\\ \textit{Options:} A) pollies; B) mollies; C) dollies; D) lollies
\\ \textit{Note:} The term "mollies" refers specifically to a subculture of homosexual men in seventeenth and eighteenth century London, known for their distinct social gatherings and practices that deviated from the dominant heterosexual norms of the time.
\item \textbf{Answer:} D
\\ \textit{Options:} A) an urban set, involved in civic bodies and voluntary associations; B) too diverse to have a strong sense of class consciousness; C) often involved in 'white collar' work; D) all of the above
\\ \textit{Note:} The correct answer "all of the above" reflects the multifaceted nature of the middle classes in the nineteenth century, encompassing their economic, social, and cultural dimensions, which collectively contributed to their distinct identity and influence in society.
\item \textbf{Answer:} C
\\ \textit{Options:} A) ethnic communities, based on shared identity and experiences of discrimination; B) gay villages, which are formed in certain parts of large cities; C) sociological communities, formed by unpopular lecturers; D) virtual communities that exist only in cyberspace
\\ \textit{Note:} Sociological communities, formed by unpopular lecturers, are established within the existing framework of academic and social structures rather than representing a novel type of community, distinguishing them from emerging forms of community that arise from new social dynamics or technological advancements.
\item \textbf{Answer:} D
\\ \textit{Options:} A) They are cultures as well as churches.; B) They practice separation of church and state.; C) They exclude women as clergy.; D) They are monotheistic.
\\ \textit{Note:} The correct answer is that these religious organizations are monotheistic because they all believe in the existence of a single, all-powerful deity, which is a fundamental characteristic of monotheistic faiths.
\item \textbf{Answer:} A
\\ \textit{Options:} A) embrace conventional cultural values but offer new means of achieving them; B) react against the loss of any meaningful religious content in the teachings of churches; C) adopt an attitude of mild disapproval towards mainstream social values; D) reject both the goals and means of conventional society and provide utopian alternatives
\\ \textit{Note:} World-affirming religions focus on integrating and enhancing societal values, providing adherents with alternative pathways to personal fulfillment and success while simultaneously affirming the importance of conventional cultural norms.
\end{enumerate}

\section*{True / False}
\begin{enumerate}[label=\arabic*.]
\setcounter{enumi}{5}
\item \textbf{Answer:} False
\\ \textit{Options:} A) True; B) False
\\ \textit{Note:} The statement is false because Giddens (1998) emphasized the importance of individualization and the changing nature of relationships in modern society rather than specifically focusing on the democratization of the family as a key aspect of New Labour's 'third way'.
\item \textbf{Answer:} True
\\ \textit{Options:} A) True; B) False
\\ \textit{Note:} Covert participant observation raises ethical concerns because it involves deceiving individuals by concealing the researcher's identity and purpose, which can undermine trust and violate principles of informed consent in research ethics.
\item \textbf{Answer:} True
\\ \textit{Options:} A) True; B) False
\\ \textit{Note:} Hegemony is understood as the way in which a dominant group maintains its power and influence by establishing and normalizing its ideology, which subsequently legitimizes and reinforces its economic, political, and cultural dominance within society.
\item \textbf{Answer:} True
\\ \textit{Options:} A) True; B) False
\\ \textit{Note:} Differential association theory posits that individuals acquire behaviors, including criminal ones, through interactions and communication with others in their social environments, thereby demonstrating that social groups significantly influence individual actions.
\item \textbf{Answer:} False
\\ \textit{Options:} A) True; B) False
\\ \textit{Note:} The statement is false because the development of transport systems in the nineteenth century, such as railways and streetcars, significantly improved access to urban areas, making urban living more appealing due to increased job opportunities and amenities compared to rural life.
\end{enumerate}

\section*{Long Form Response}
\begin{enumerate}[label=\arabic*.]
\setcounter{enumi}{10}
\item \textbf{Answer:} The United States government defines poverty through the federal poverty line, which was established in the 1960 s by sociologist Molly Orshansky. This measure was originally based on the cost of a minimum food diet multiplied by three, reflecting the economic conditions of that era. Since then, the poverty line has been adjusted for inflation, allowing it to better represent the current economic realities faced by households. These adjustments ensure that the poverty threshold remains relevant, yet criticisms have arisen regarding its accuracy in capturing the complexities of poverty today, such as regional cost-of-living variations and non-food expenses. Overall, while the poverty line serves as a critical tool for policy-making and resource allocation, it is rooted in a historical context that may not fully encompass the multifaceted nature of poverty in contemporary society.
\\ \textit{Note:} correct because it accurately describes the federal poverty line's origin based on a food cost analysis, its establishment by Molly Orshansky in the 1960 s, and acknowledges the necessary adjustments for inflation to reflect contemporary economic conditions.
\item \textbf{Answer:} Booth's 1901 study on poverty in London is significant as it revealed that approximately 30.70\% of the population was living in poverty. This finding highlighted the stark social inequalities present in urban environments during that era, illustrating the extent to which economic hardship affected a substantial portion of the populace. By categorizing individuals into various classes based on their living conditions, Booth's research provided a framework for understanding the systemic nature of poverty and its intersection with social stratification. His analysis not only raised awareness about the plight of the impoverished but also prompted discussions around social reform and the role of government in addressing inequality. Ultimately, Booth's work laid the groundwork for future sociological studies on poverty and social justice.
\\ \textit{Note:} correct because it accurately identifies the significant percentage of the population living in poverty as revealed by Booth's study, emphasizing the implications of this finding for understanding the systemic social inequalities prevalent in early 20 th-century urban society.
\end{enumerate}

\vfill
\noindent\textit{End of Answer Key}
\end{document}